% TEMPLATE-MILC-v3.tex - plantilla maestra UNICA de las guias MILC v3.
%
% La usan las cuatro familias de guias, cada una con su driver:
%   · Grados 6-11          -> scripts/build-guias-g11.py
%   · Bebras               -> scripts/build-guias-bebras.py
%   · Semillero            -> scripts/build-guias-semillero.py
%   · Territorio interior  -> scripts/build-guias-territorio-interior.py
%
% Antes habia tres copias casi identicas (~400 lineas iguales, 6 distintas):
% cada mejora habia que aplicarla tres veces, y en la practica no se hacia
% (el motor de recursos vivia solo en dos de ellas). Lo que varia entre
% familias esta parametrizado en 6 placeholders que cada driver llena
% (se nombran aqui SIN su sintaxis real: el builder reemplaza por texto plano,
% tambien dentro de los comentarios, y un valor multilinea rompe el preambulo):
%
%   HEADER_IZQ        encabezado izquierdo de pagina
%   HEADER_DER        encabezado derecho de pagina
%   PORTADA_ETIQUETA  rotulo sobre el numero grande (GRADO/MOMENTO/MODULO)
%   PORTADA_SERIE     linea de serie de la portada
%   PORTADA_ID        identificador de la guia en la portada
%   PORTADA_CIRCULO   el circulito de grado (°), o vacio si no aplica
%   PDF_TITULO        titulo en texto plano (sin \textbf ni comillas TeX) para
%                     los metadatos del PDF (/Title); lo produce lib_xelatex.texto_pdf
%
% Composicion (auditoria 2026-09-03): las bandas de estacion piden espacio
% (\Needspace*) y prohiben el salto justo despues (\nopagebreak), asi no quedan
% huerfanas al pie; las cajas largas (softbox, drawbox, infoband, puentebox) son
% `breakable`, asi una caja de media pagina ya no arrastra todo a la siguiente
% dejando la anterior al 40 %. Todo texto sobre mostaza o verde va en milcNegro
% (blanco sobre #E5B400 da 1,93:1 y sobre #58A923 2,95:1; ninguno pasa WCAG AA).
% El resto de claves las documenta content/guias/_SCHEMA.md. El script valida
% que no queden placeholders sin reemplazar antes de escribir el archivo final.

% Etiquetado PDF (latex-lab): /Lang, estructura y texto alternativo de las
% figuras, para lectores de pantalla. Debe ir ANTES de \documentclass.
\DocumentMetadata{lang=es-CO,tagging=on}
\documentclass[11pt,letterpaper]{article}
% headheight=24pt: la cabecera derecha lleva dos lineas (identidad de la guia
% y estacion actual). headsep se acorta en la misma medida para que el bloque
% de texto quede exactamente donde estaba con headheight=12pt.
\usepackage[margin=2.0cm,headheight=24pt,headsep=13pt]{geometry}
\usepackage[table]{xcolor}
\usepackage{fontspec}
\usepackage[spanish]{babel}
\usepackage{tikz}
% Librerías TikZ para los diagramas de las guías (bloque `recursos:`):
%  · babel  → babel en español vuelve activos caracteres como «>», y TikZ los usa
%             en sus opciones (>=stealth, ->). Sin esta librería, cualquier
%             diagrama con flechas revienta con «\language@active@arg> has an
%             extra }». Debe ir ANTES de que se dibuje nada.
%  · positioning        → sintaxis «below left=2pt and 3pt».
%  · decorations.pathreplacing → llaves ({brace}) para señalar rangos.
\usetikzlibrary{babel,positioning,decorations.pathreplacing}
\usepackage{graphicx}
% Circuitos: el trabajo de electrónica se hace hoy con micro:bit y sus sensores
% integrados (MakeCode, sin cableado), así que no se necesitan esquemáticos.
% Si algún día entran montajes con protoboard, basta descomentar esta línea:
% los diagramas .tex/.tikz declarados en `recursos:` ya se incrustan solos.
% \usepackage{circuitikz}
\usepackage[most]{tcolorbox}
\usepackage{tabularx}
\usepackage{array}
\usepackage{fancyhdr}
\usepackage{titlesec}
\usepackage[document]{ragged2e}  % bandera a la derecha en todo el cuerpo
\usepackage{enumitem}
\usepackage{microtype}
\usepackage{etoolbox}
\usepackage{needspace}
% hyperref al final del preambulo: metadatos del PDF (titulo, autor, idioma,
% familia), marcadores por estacion (\pdfbookmark) y enlaces sin recuadro.
\usepackage[hidelinks,
  pdftitle={Tablas dinámicas — agrupar y comparar},
  pdfauthor={PhD. Álvaro Cárdenas Orozco},
  pdfsubject={Tecnología e Informática · Grado 9 · Guía 25 · MILC v3},
  pdflang={es-CO},
  bookmarksopen=true,bookmarksnumbered=false]{hyperref}

% Helvetica Neue no trae la flecha U+2192, asi que cada «→» escrito literalmente
% en el YAML se compilaba a NADA, en silencio (462 flechas en 61 guias: rutas del
% tipo «paleta → tipografia → lockup» salian rotas). Mapeamos el caracter a su
% version matematica, que si tiene glifo. Asi el autor puede seguir escribiendo
% «→» comodamente en el YAML.
\usepackage{newunicodechar}
\newunicodechar{→}{\ensuremath{\rightarrow}}
% Mismo problema con los simbolos de marca y vinetas: el «✓» de «una palomita ✓»
% salia como cuadrito vacio en el PDF de todas las guias que lo usaban.
\usepackage{pifont}
\newunicodechar{✓}{\ding{51}}
\newunicodechar{✗}{\ding{55}}
\newunicodechar{✔}{\ding{52}}
\newunicodechar{•}{\textbullet}
\newunicodechar{≥}{\ensuremath{\geq}}
\newunicodechar{≤}{\ensuremath{\leq}}

\setmainfont{Helvetica Neue}
\setsansfont{Helvetica Neue}
\setlength{\parindent}{0pt}
\setlength{\parskip}{6pt}
% Sin particion de palabras: en 6.º-7.º una palabra cortada por guion al final
% de linea («Comuni-cacion») frena la lectura mucho mas que una linea corta.
\hyphenpenalty=10000
\exhyphenpenalty=10000
\pretolerance=2000
\tolerance=3000
\setlength{\emergencystretch}{3em}
\linespread{1.10}
\renewcommand{\arraystretch}{1.25}
\setlist{leftmargin=5mm,itemsep=1mm,topsep=1mm}
\newcolumntype{Y}{>{\RaggedRight\arraybackslash}X}

\definecolor{milcMagenta}{HTML}{E6007E}
\definecolor{milcVino}{HTML}{5A0038}
\definecolor{milcTurquesa}{HTML}{0093A5}
\definecolor{milcVerde}{HTML}{58A923}
\definecolor{milcMostaza}{HTML}{E5B400}
\definecolor{milcOcre}{HTML}{B86600}
\definecolor{milcNegro}{HTML}{241816}
\definecolor{milcGris}{HTML}{F7F7F4}
\definecolor{milcLinea}{HTML}{E8E2DD}
\definecolor{milcGradePrimary}{HTML}{7C3AED}
\definecolor{milcGradeDark}{HTML}{4C1D95}
\definecolor{milcGradeSoft}{HTML}{EDE9FE}
\definecolor{milcGradeAccent}{HTML}{CFE7FF}
\definecolor{milcGradeText}{HTML}{FFFFFF}
\color{milcNegro}

\pagestyle{fancy}
\fancyhf{}
% Cabecera de dos lineas: identidad de la guia arriba y, a la derecha y en
% gris, la estacion en curso (\markright desde \milcestacion). Antes de la
% primera estacion la segunda linea va vacia; el \strut mantiene la altura
% para que la linea de arriba no baile entre paginas.
\lhead{\small Tecnología e Informática · Grado 9\\[-1pt]{\footnotesize\strut}}
\rhead{\small Guía 25 · MILC v3\\[-1pt]{\footnotesize\strut\color{milcNegro!65}\rightmark}}
\cfoot{\small\thepage}
\renewcommand{\headrulewidth}{0pt}

% Sobre mostaza (#E5B400) y verde (#58A923) el texto va en milcNegro; sobre el
% resto de la paleta, en blanco. Se usa dentro de fonttitle/fontupper, donde un
% \color posterior manda sobre coltitle.
\newcommand{\milctintasobre}[1]{%
  \ifboolexpr{test{\ifstrequal{#1}{milcMostaza}} or test{\ifstrequal{#1}{milcVerde}}}%
    {\color{milcNegro}}{\color{white}}}

\newtcolorbox{titlebox}[1]{
  enhanced,colback=#1,colframe=#1,arc=7mm,boxrule=0pt,
  left=7mm,right=7mm,top=3mm,bottom=3mm,
  before skip=8pt,after skip=8pt,
  fontupper=\sffamily\bfseries\Large\color{white}
}

\newtcolorbox{softbox}[3]{
  enhanced,breakable,pad at break=3mm,
  colback=#2,colframe=#1,borderline west={4pt}{0pt}{#1},
  arc=2mm,boxrule=.4pt,left=4mm,right=4mm,top=3mm,bottom=3mm,
  before skip=6pt,after skip=8pt,title={#3},coltitle=white,
  colbacktitle=#1,fonttitle=\sffamily\bfseries\milctintasobre{#1},fontupper=\RaggedRight,
  attach boxed title to top left={xshift=3mm,yshift=-2mm},
  boxed title style={arc=2mm, boxrule=0pt}
}

\newtcolorbox{drawbox}[2]{
  enhanced,breakable,pad at break=4mm,
  colback=white,colframe=#1,arc=4mm,boxrule=1pt,
  left=5mm,right=5mm,top=4mm,bottom=4mm,before skip=8pt,after skip=8pt,
  title={#2},coltitle=white,colbacktitle=#1,fonttitle=\sffamily\bfseries\milctintasobre{#1},
  fontupper=\RaggedRight,
  attach boxed title to top left={xshift=4mm,yshift=-2mm},
  boxed title style={arc=3mm, boxrule=0pt}
}

\newtcolorbox{infoband}[1]{
  enhanced,breakable,pad at break=2mm,colback=#1!8,colframe=#1!50,arc=2mm,boxrule=.5pt,
  left=4mm,right=4mm,top=2mm,bottom=2mm,before skip=4pt,after skip=4pt,
  fontupper=\small\RaggedRight
}

% Puente narrativo entre secciones (contrato editorial Conéctate).
% Da continuidad: "Acabas de X. Ahora vas a Y porque Z."
% Visualmente: banda izquierda en color del grado, texto en itálica pequeña.
\newtcolorbox{puentebox}{
  enhanced,breakable,pad at break=2mm,colback=milcGradeSoft,colframe=milcGradeDark,
  borderline west={3pt}{0pt}{milcGradeDark},
  arc=0mm,boxrule=0pt,left=5mm,right=4mm,top=2.5mm,bottom=2.5mm,
  before skip=4pt,after skip=8pt,
  fontupper=\itshape\small\color{milcNegro}\RaggedRight
}
% La flecha va en modo matematico a proposito: Helvetica Neue NO tiene el glifo
% U+2192, asi que \textrightarrow se compilaba a NADA («Missing character: There
% is no → in font Helvetica Neue») y el puente salia sin flecha. En modo
% matematico la flecha viene de la fuente matematica y si se dibuja.
\newcommand{\puente}[1]{%
  \begin{puentebox}%
    \textcolor{milcGradeDark}{$\rightarrow$}\hspace{2mm}#1%
  \end{puentebox}%
}

\newcommand{\linead}[1]{\noindent\color{milcLinea}\rule{#1}{0.5pt}\color{milcNegro}}
\newcommand{\respuesta}[1]{\par\vspace{2mm}\foreach \n in {1,...,#1}{\noindent\color{milcLinea}\rule{\linewidth}{0.5pt}\par\vspace{5mm}}\color{milcNegro}}
\newcommand{\checkbox}{\textcolor{milcGradeDark}{\fbox{\phantom{x}}}\hspace{5pt}}

% ─── Listas aireadas para las guías socioemocionales ────────────────────
% Componer las verificaciones como «\checkbox texto\\» deja el texto que salta
% de línea debajo del cuadrito y apelmaza el bloque. Estos entornos alinean la
% continuación y airean los ítems. Los usa Territorio Interior; cualquier
% familia puede hacerlo.
\newlist{checklista}{itemize}{1}
\setlist[checklista]{
  label=\textcolor{milcGradeDark}{\fbox{\phantom{x}}},
  leftmargin=9mm, labelsep=3.2mm, itemsep=2.4mm,
  topsep=2.5mm, parsep=0pt, partopsep=0pt, before=\RaggedRight
}
\newlist{pasoslista}{enumerate}{1}
\setlist[pasoslista]{
  label=\textcolor{milcGradeDark}{\sffamily\bfseries\arabic*.},
  leftmargin=8mm, labelsep=3mm, itemsep=2.8mm,
  topsep=1mm, parsep=0pt, partopsep=0pt, before=\RaggedRight
}

\newcommand{\phasecard}[4]{
\begin{tcolorbox}[enhanced,colback=#3,colframe=#2,arc=3mm,boxrule=.5pt,height=3.05cm,valign=center,halign=center,left=1.5mm,right=1.5mm,top=2mm,bottom=2mm]
{\sffamily\bfseries\fontsize{22}{24}\selectfont\textcolor{#2}{#1}}\par
{\sffamily\bfseries\small #4}
\end{tcolorbox}
}

% ── Piezas del rediseno 2026 (legibilidad 6.º-11.º) ───────────────────
% Banda de estacion: reemplaza el titlebox macizo. Titulo a la izquierda,
% dato practico (tiempo/modalidad) a la derecha, en una sola linea.
% Nunca huerfana: pide 9 lineas libres (o salta de pagina rellenando con
% \vfill) y prohibe el salto justo despues de la banda. Nueve y no seis:
% la caja partible que sigue necesita ~80pt (titulo adosado, relleno y dos
% lineas) para arrancar en la misma pagina; con menos, tcolorbox la manda
% entera a la siguiente y la banda se queda sola. El penalty 10000 va
% en `after`, ANTES del espacio posterior: un `after skip` (aunque sea 0pt)
% es un \vskip tras la caja, y ese glue es un punto de corte legal que un
% \nopagebreak escrito despues de \end{tcolorbox} ya no cierra. Cada banda
% deja marcador PDF y actualiza la cabecera derecha.
\newcounter{milcestacion}
\newcommand{\milcestacion}[2]{%
\Needspace*{9\baselineskip}%
\stepcounter{milcestacion}%
\pdfbookmark[1]{#2}{milcestacion\themilcestacion}%
\markright{#2}%
\begin{tcolorbox}[enhanced,colback=#1,colframe=#1,arc=2mm,boxrule=0pt,
  left=5mm,right=5mm,top=2.6mm,bottom=2.6mm,before skip=11pt,
  after={\par\nopagebreak\vskip7pt}]
{\sffamily\bfseries\fontsize{15}{18}\selectfont\milctintasobre{#1}#2}
\end{tcolorbox}}

% Tarjeta de paso numerada: sustituye las tablas «Pilar / Aplicacion», que
% eran formato de consulta docente, no de estudiante. El numero va en un
% circulo sobre el borde para que la secuencia se lea de un vistazo.
\newcommand{\milcpaso}[3]{%
\Needspace{4\baselineskip}%
\begin{tcolorbox}[enhanced,colback=white,colframe=milcLinea,arc=1.5mm,boxrule=.9pt,
  left=14mm,right=4mm,top=3mm,bottom=3mm,before skip=4pt,after skip=4pt,
  fontupper=\RaggedRight,
  overlay={\node[anchor=north west,circle,fill=#2,text=white,inner sep=0pt,
    minimum size=7.4mm,font=\sffamily\bfseries\small]
    at ([xshift=3.6mm,yshift=-3.2mm]frame.north west) {#1};}]
#3
\end{tcolorbox}}

% Casilla marcable de tamano util con lapiz (la anterior era \fbox{\phantom{x}},
% de alto variable segun el texto que la seguia).
\newcommand{\milcmarca}{\framebox[3.8mm]{\rule{0pt}{3.4mm}}}

% Fila marcable de lista suelta (checks de la estacion 1).
\newcommand{\milccheckrow}[1]{%
\par\vspace{1.8mm}\noindent
\makebox[6.4mm][l]{\raisebox{-.55mm}{\textcolor{milcNegro}{\milcmarca}}}%
\parbox[t]{\dimexpr\linewidth-6.4mm\relax}{\RaggedRight #1}\par}

% Celda marcable dentro de la rubrica (tabla de autoevaluacion). Misma
% sangria francesa, pero pensada para vivir dentro de una columna X.
\newcommand{\milcopcion}[1]{%
\makebox[5.2mm][l]{\raisebox{-.55mm}{\textcolor{milcNegro}{\milcmarca}}}%
\parbox[t]{\dimexpr\linewidth-5.2mm\relax}{\RaggedRight #1}}

% ── Recursos visuales de la guía ──────────────────────────────────────
% Declarados en el bloque `recursos:` del YAML y emitidos por el builder.
% \guiaFigura   → imagen rasterizada o PDF (foto, carta celeste, montaje).
% \guiaDiagrama → fuente TikZ/circuitikz (.tex) incrustada como vector:
%                 es la vía para circuitos y esquemáticos editables.
% [#1] = texto alternativo (alt) · #2 = ruta relativa al .tex · #3 = pie
% El alt llega al PDF etiquetado (/Alt) y lo lee el lector de pantalla.
\newcommand{\guiaFigura}[3][]{%
\begin{center}
\includegraphics[width=0.88\linewidth,height=0.38\textheight,keepaspectratio,alt={#1}]{#2}\\[1.5mm]
{\footnotesize\itshape\textcolor{milcNegro!75}{#3}}
\end{center}
}

\newcommand{\guiaDiagrama}[2]{%
\begin{center}
\input{#1}\\[1.5mm]
{\footnotesize\itshape\textcolor{milcNegro!75}{#2}}
\end{center}
}

\begin{document}
\thispagestyle{empty}

% ── Portada ───────────────────────────────────────────────────────────
% Antes: un rectangulo de color a pagina completa. Se veia bien en pantalla,
% pero en la fotocopiadora del colegio se comia el toner y salia gris sucio.
% Ahora la banda ocupa el tercio superior y el resto va en blanco.
\begin{tikzpicture}[remember picture,overlay]
  \path[fill=milcGradePrimary]
    (current page.north west) rectangle ([yshift=-10.6cm]current page.north east);

  \node[anchor=north west,text=milcGradeText] at ([xshift=2.0cm,yshift=-1.55cm]current page.north west)
    {\sffamily\bfseries\fontsize{12}{14}\selectfont GRADO};
  \node[anchor=north west,text=milcGradeText,opacity=.85] at ([xshift=2.0cm,yshift=-2.35cm]current page.north west)
    {\sffamily\bfseries\fontsize{8.5}{10}\selectfont SERIE GUÍAS · TECNOLOGÍA E INFORMÁTICA};
  \node[anchor=north east,text=milcGradeText,opacity=.85,align=right] at ([xshift=-2.0cm,yshift=-1.55cm]current page.north east)
    {\sffamily\bfseries\fontsize{9}{12}\selectfont I.E. Sor María Juliana\\Cartago, Valle del Cauca};

  \node[anchor=west,text=milcGradeText] (gradonum) at ([xshift=1.85cm,yshift=-7.1cm]current page.north west)
    {\sffamily\bfseries\fontsize{112}{112}\selectfont 9};
  % El circulito de grado (°) solo aplica a las guias de grado («11°»); en
  % Bebras y Semillero el numero grande es un momento/modulo, no un grado.
  % Se posiciona relativo al nodo `gradonum`, no en coordenadas absolutas.
  \draw[milcGradeText,line width=1.6pt]
    ([xshift=2.0mm,yshift=-4.5mm]gradonum.north east) circle (.15cm);

  \node[anchor=west,text=milcGradeText,text width=11.4cm,align=left] at ([xshift=1.5cm,yshift=0.5cm]gradonum.east)
    {
     {\sffamily\bfseries\fontsize{9}{11}\selectfont\MakeUppercase{Período 3 · Datos --- del registro al insight}}\\[3mm]
     {\sffamily\bfseries\fontsize{21}{25}\selectfont\setlength{\baselineskip}{27pt}Tablas dinámicas ---\\[4pt]agrupar y\\[4pt]comparar}\\[3.5mm]
     {\sffamily\bfseries\fontsize{15}{18}\selectfont Guía 25-9-TIC}
    };
\end{tikzpicture}

\vspace*{9.4cm}
\vfill

{\sffamily\bfseries\fontsize{8}{10}\selectfont\textcolor{milcGradeDark}{LO QUE TE LLEVAS HOY}}

\begin{tcolorbox}[enhanced,colback=white,colframe=milcNegro,arc=2mm,boxrule=1.6pt,
  left=5mm,right=5mm,top=4mm,bottom=4mm,before skip=3pt,after skip=12pt,
  fontupper=\RaggedRight\fontsize{11}{15.5}\selectfont]
1 tabla dinámica que responda una pregunta agrupando los datos, con captura y conclusión de 3 renglones
\end{tcolorbox}

\begin{tcolorbox}[enhanced,colback=milcGris,colframe=milcGris,arc=2mm,boxrule=0pt,
  left=5mm,right=5mm,top=3.5mm,bottom=3.5mm,after skip=12pt]
{\sffamily\bfseries\fontsize{8}{10}\selectfont\textcolor{milcNegro!55}{ESTA GUÍA ES DE}}\\[3mm]
\begin{tabularx}{\linewidth}{X p{3.2cm} p{3.2cm}}
\linead{\linewidth} & \linead{3.0cm} & \linead{3.0cm}\\[-1mm]
{\fontsize{8.5}{10}\selectfont\textcolor{milcNegro!55}{Nombre completo}} &
{\fontsize{8.5}{10}\selectfont\textcolor{milcNegro!55}{Grupo}} &
{\fontsize{8.5}{10}\selectfont\textcolor{milcNegro!55}{Fecha}}\\
\end{tabularx}
\end{tcolorbox}

\vfill

\noindent\textcolor{milcLinea}{\rule{\linewidth}{1pt}}\\[2mm]
{\sffamily\bfseries\fontsize{9.5}{12}\selectfont Autor: Dr. Álvaro Cárdenas Orozco}\\[1mm]
{\sffamily\fontsize{8.5}{11}\selectfont\textcolor{milcNegro!55}{Metodología MILC · apertura ancestral · pensamiento computacional · triángulo Dussel–estoicismo–Floridi}}

\newpage

\section*{Apertura: Tablas dinámicas --- agrupar y comparar}

\begin{softbox}{milcGradeDark}{milcGradeSoft}{Saber ancestral}
En el Eje Cafetero, la palabra «cuadrilla» nombra dos grupos que no se parecen en nada. Uno es el de los recolectores que entran al cafetal detrás de un patrón de corte. El patrón marca los surcos con una bandera y reparte el trabajo (Saldarriaga Ramírez, 2024). Ahí se pertenece mientras dure la cosecha. Y a la gente se la agrupa por de dónde viene: paisa, tolimense, pastuso. El otro son las Cuadrillas del Carnaval de Riosucio. Doce a quince personas trabajan un año entero en una letra, una música y un disfraz (Ministerio de Cultura, 2011). Ninguna de las tres cosas se puede repetir nunca más. Las dos cuadrillas dan pertenencia. Pero una la da mientras uno sirva, y la otra mientras uno cree. La \textbf{cara de exclusión} está del lado del cafetal. Se trabaja sin contrato ni seguridad social, en cuarteles hacinados, con jornadas que arrancan a las 6:30 de la mañana. Y hacia afuera pesa el estigma de «mal vestidos» o «bebedores de bar». Pertenecer a una cuadrilla puede convivir con quedar excluido del pueblo entero.\par\smallskip{\footnotesize\textcolor{milcNegro!70}{\textbf{Fuente.} Saldarriaga Ramírez, C. (2024). Andariegos: prácticas culturales de los recolectores itinerantes de café del municipio de Pereira, Colombia. \emph{Apuntes: Revista de Estudios sobre Patrimonio Cultural, 37}. https://doi.org/10.11144/Javeriana.apu37.apcr}}
\end{softbox}

\begin{softbox}{milcTurquesa}{milcGris}{Saber contemporáneo}
Una \textbf{tabla dinámica} toma una tabla larga y la resume por grupos. Agrupa por una columna y calcula un indicador sobre otra. La sesión 3 te dio una suma total; esta te da \emph{suma por categoría}. Total por mes, promedio por curso, conteo por barrio. Responde a tres decisiones simples: qué agrupo (Filas), qué calculo (Valores), qué cruzo (Columnas). Sin programar nada contesta preguntas que a mano tardarían horas. Y hace visible lo que un promedio general esconde.
\end{softbox}

\begin{softbox}{milcMostaza}{milcGris}{Pregunta-puente}
\textit{¿Por qué la misma palabra nombra dos cuadrillas que no se parecen en nada? ¿Y qué se pierde cuando un promedio mete en una sola fila a gente que no vive lo mismo?}
\end{softbox}

\begin{softbox}{milcVerde}{milcGris}{Saber y saber hacer}
Al terminar podrás: (1) \textbf{analizar} cuándo una pregunta necesita agruparse y cuándo basta una fórmula simple; (2) \textbf{explicar} las cuatro zonas de una tabla dinámica y qué operación pide cada pregunta; (3) \textbf{crear} tu propia tabla dinámica con una conclusión que nombre el grupo destacado y el grupo en riesgo.
\end{softbox}



\small
\begin{tabularx}{\textwidth}{>{\bfseries}p{3.0cm}Y}
\rowcolor{milcGradePrimary!12}Autor & Dr. Álvaro Cárdenas Orozco\\
DBA & Tratamiento y visualización honesta de datos para la toma de decisiones (MEN, T\&I)\\
Referentes & Ética del dato · Visualización honesta · Pensamiento computacional\\
Producto final & 1 tabla dinámica que responda una pregunta agrupando los datos, con captura y conclusión de 3 renglones\\
\end{tabularx}
\normalsize

\puente{La sesión 3 te dio fórmulas para toda la tabla. La sesión 4 te dio filtros para cortarla. Hoy aprendes a \emph{re-agruparla}. Ya no «cuánto en total», sino «cuánto por categoría». Es el salto del cálculo único al cálculo comparable.}

\milcestacion{milcGradeDark}{Ruta MILC: así vamos a aprender}
\begin{tabularx}{\textwidth}{>{\centering\arraybackslash}X>{\centering\arraybackslash}X>{\centering\arraybackslash}X>{\centering\arraybackslash}X}
\phasecard{1}{milcVerde}{milcGris}{Escucha\\[-1mm]\small Observo, escucho y pregunto.} &
\phasecard{2}{milcTurquesa}{milcGris}{Entender\\[-1mm]\small Sistematizo y ordeno.} &
\phasecard{3}{milcMagenta}{milcGris}{Hacer\\[-1mm]\small Construyo y aplico.} &
\phasecard{4}{milcVino}{milcGris}{Cierre\\[-1mm]\small Evalúo y reflexiono.}
\end{tabularx}


\puente{Antes de teorizar, vas a contestar \emph{a ojo} dos preguntas de agrupación sobre tu tabla. Después las vas a ver resueltas en dos clics. El contraste te muestra para qué sirve de verdad la herramienta.}

\milcestacion{milcVerde}{Estación 1 de 3 · Escucha}

\begin{softbox}{milcVerde}{milcGris}{Escena para escuchar}
\textbf{Actividad 1 · ANALIZA --- ¿Total o por grupo?} (15 min · individual). (1) Toma tu tabla de la sesión 2 con las fórmulas de la sesión 3 ya aplicadas. (2) Responde \emph{a ojo} dos preguntas: cuál categoría tiene el promedio más alto en una columna numérica, y cuál categoría tiene más filas. (3) Cronometra cuánto tardas y anota si dudaste. (4) Abre Insertar $\to$ Tabla dinámica, arrastra esa categoría a \emph{Filas} y la columna numérica a \emph{Valores}. (5) Verifica las dos respuestas y compara qué método fue más rápido y cuál dio el resultado correcto.
\end{softbox}

{\sffamily\bfseries\fontsize{8}{10}\selectfont\textcolor{milcNegro!55}{MARCA CUANDO LO HAYAS HECHO}}
\milccheckrow{Respondí a ojo cuál categoría destaca por promedio y cuál por conteo.}
\milccheckrow{Construí una primera tabla dinámica para verificar las dos respuestas.}
\milccheckrow{Comparé honestamente: cuál método fue más rápido y cuál dio el resultado correcto.}

\begin{infoband}{milcVerde}


\textbf{Lo que va al cuaderno (Actividad 1 · ANALIZA).} \textbf{Título:} «Actividad 1 --- ¿Total o por grupo?». \textbf{Formato:} tabla de dos filas y cuatro columnas (pregunta / respuesta a ojo / respuesta con tabla dinámica / tiempo a ojo), más una línea sobre quién ganó en velocidad y quién en precisión. \textbf{Extensión:} un tercio de página. \textbf{Sabes que terminaste cuando} reconoces en qué casos agrupar a ojo basta y en cuáles la tabla dinámica es indispensable.
\end{infoband}

\puente{Ya viste el contraste. Ahora vas a entender la \textbf{anatomía} de una tabla dinámica. Las cuatro zonas, el tipo de operación y la diferencia entre agrupar por categoría y agrupar por fecha.}

\milcestacion{milcTurquesa}{Estación 2 de 3 · Entender}

\begin{softbox}{milcTurquesa}{milcGris}{Lo que vas a entender}
\textbf{Actividad 2 · EXPLICA --- Las cuatro zonas} (20 min · parejas). (1) Con tu pareja, escriban las cuatro zonas de una tabla dinámica con la función de cada una. (2) Escriban la secuencia de tres pasos que sigue toda tabla dinámica. (3) Tomen tres preguntas de sus tablas y digan qué operación necesita cada una. (4) Busquen una pregunta que solo se responda cruzando dos categorías y escríbanla completa.
\end{softbox}

\milcpaso{1}{milcTurquesa}{La tabla dinámica se descompone en cuatro zonas. \textbf{Filas} es la categoría por la que agrupas: mes, curso, barrio, área. \textbf{Valores} es la columna numérica sobre la que calculas. \textbf{Columnas} es opcional y sirve para cruzar una segunda categoría. \textbf{Filtros} también es opcional y limita toda la tabla. Sin Filas y Valores no hay tabla dinámica.}
\milcpaso{2}{milcTurquesa}{Toda tabla dinámica sigue el patrón «categoría $\to$ medida $\to$ operación». Eliges qué agrupa, qué se calcula y con qué operación: SUMA, PROMEDIO, CUENTA, MAX o MIN. El error más común es confundir CUENTA con SUMA. CUENTA dice cuántos hay; SUMA dice cuánto suman. Una columna de notas pide PROMEDIO.}
\milcpaso{3}{milcTurquesa}{Todo se juega en una pregunta: ¿qué decisión cambia al ver los grupos por separado? Un promedio general de 6.2 parece regular. Pero si 9A da 7.5 y 9B da 4.9, ahí está el problema. La tabla dinámica sirve cuando el total esconde diferencias que importan para decidir.}
\milcpaso{4}{milcTurquesa}{Algoritmo para construir una: (1) escribe la pregunta de agrupación; (2) selecciona toda la tabla con encabezados; (3) Insertar $\to$ Tabla dinámica $\to$ Nueva hoja; (4) arrastra la categoría a Filas; (5) arrastra la columna numérica a Valores; (6) verifica la operación con clic derecho; (7) interpreta y captura.}

\begin{drawbox}{milcTurquesa}{Anatomía de una tabla dinámica bien construida}
\textbf{Pregunta:} qué comparas y por qué. \\[1mm] \textbf{Filas:} la categoría que agrupa. \\[1mm] \textbf{Valores:} la columna numérica que calculas. \\[1mm] \textbf{Operación:} SUMA / PROMEDIO / CUENTA / MAX / MIN. \\[1mm] \textbf{Resultado:} qué categoría destaca y cuál queda al final. \\[1mm] \textbf{Interpretación:} una frase que conecta el resultado con una decisión posible.
\end{drawbox}

\begin{softbox}{milcMostaza}{milcGris}{Errores comunes y su corrección}
(1) Arrastrar a Valores una columna de texto: Excel cuenta filas en vez de calcular. (2) Olvidar cambiar la operación, que arranca en SUMA aunque quieras PROMEDIO. (3) Agrupar por una columna con datos sucios, donde «9A» y «9a» aparecen como dos grupos. (4) Construir la tabla sin pregunta escrita y terminar con un resumen lindo sin conclusión. (5) Reportar la tabla sin nombrar la categoría que destaca.
\end{softbox}

\begin{infoband}{milcTurquesa}


\textbf{Lo que va al cuaderno (Actividad 2 · EXPLICA).} \textbf{Título:} «Actividad 2 --- Las cuatro zonas». \textbf{Formato:} las cuatro zonas con su función, la secuencia de tres pasos, las tres preguntas con su operación y la pregunta que necesita cruzar dos categorías. \textbf{Extensión:} media página. \textbf{Sabes que terminaste cuando} distingues una operación técnicamente correcta de una operación que responde tu pregunta real.
\end{infoband}

\puente{Tienes el método. Toca aplicarlo. Una sola tabla dinámica bien hecha, con pregunta escrita, agrupación clara, cálculo correcto y conclusión de tres renglones.}

\milcestacion{milcMagenta}{Estación 3 de 3 · Hacer}

\begin{softbox}{milcMagenta}{milcGris}{Cómo vas a construir tu producto}
\textbf{Actividad 3 · CREA --- Tu tabla dinámica con conclusión} (30 min · individual). (1) Escribe en una frase la pregunta de agrupación que quieres responder. (2) Elige la categoría que va a Filas y la columna numérica que va a Valores. (3) Construye la tabla dinámica y ajusta la operación. (4) Ordena los resultados de mayor a menor. (5) Toma una captura con la tabla visible. (6) Escribe tres renglones que nombren el grupo destacado, el grupo en riesgo y la decisión posible.
\end{softbox}

\milcpaso{1}{milcMagenta}{Tu producto se descompone en cuatro piezas: la pregunta escrita en una frase, la tabla dinámica configurada, la captura de pantalla y la conclusión de tres renglones que nombre el grupo destacado y el grupo en riesgo.}
\milcpaso{2}{milcMagenta}{Sigue el patrón «categoría con sentido pedagógico, social o económico». No agrupes por columnas técnicas como el ID o la fecha exacta. Agrupa por categorías que produzcan decisión: curso, área, mes, barrio, tipo de gasto. Una buena tabla dinámica genera movimiento; una mala genera decoración.}
\milcpaso{3}{milcMagenta}{Cada conclusión sigue un guion mínimo: el grupo X destaca porque…, el grupo Y queda atrás porque…, esto sugiere que… Tres frases cortas. Si tu conclusión ocupa más de tres renglones, recórtala.}
\milcpaso{4}{milcMagenta}{Algoritmo: (1) escribe la pregunta; (2) elige la categoría; (3) elige la columna numérica; (4) construye la tabla dinámica; (5) ajusta la operación; (6) ordena de mayor a menor; (7) captura; (8) escribe los tres renglones.}

\begin{drawbox}{milcMagenta}{Checklist antes de cerrar la tabla dinámica}
\checkbox Pregunta escrita en una frase antes de abrir la tabla dinámica. \\[1mm] \checkbox Tabla dinámica con Filas y Valores configurados y la operación correcta. \\[1mm] \checkbox Resultados ordenados de mayor a menor. \\[1mm] \checkbox Captura de pantalla con la tabla dinámica visible. \\[1mm] \checkbox Conclusión de tres renglones con grupo destacado, grupo en riesgo y decisión posible.
\end{drawbox}

\begin{softbox}{milcMagenta}{milcGris}{Plantilla de trabajo}
\textbf{Pregunta.} \textit{¿Qué quieres comparar entre grupos?} \\[1mm] \textbf{Filas.} \textit{Categoría que agrupa.} \\[1mm] \textbf{Valores.} \textit{Columna numérica.} \\[1mm] \textbf{Operación.} \textit{SUMA / PROMEDIO / CUENTA / MAX / MIN.} \\[1mm] \textbf{Resultado.} \textit{Grupo destacado y grupo al final.} \\[1mm] \textbf{Conclusión.} \textit{Tres renglones: destaca / queda atrás / sugiere.}
\end{softbox}

\begin{infoband}{milcMagenta}


\textbf{Lo que va al cuaderno (Actividad 3 · CREA).} \textbf{Título:} «Actividad 3 --- Mi tabla dinámica». \textbf{Formato:} la pregunta, la configuración de zonas y operación, la captura pegada o descrita y la conclusión de tres renglones. \textbf{Extensión:} media página. \textbf{Sabes que terminaste cuando} puedes defender en voz alta qué decisión propones a partir de la tabla.
\end{infoband}

\puente{Lo que sigue resume \emph{qué} entregas hoy: una tabla dinámica, una captura y una conclusión escrita. El criterio principal es que la conclusión \textbf{nombre} qué categoría destaca y qué decisión podría seguir de ahí.}

\milcestacion{milcMostaza}{Producto de la sesión}
\textbf{1 tabla dinámica con pregunta, captura y conclusión de 3 renglones}.
\begin{softbox}{milcMostaza}{milcGris}{Criterios de calidad}
(1) Pregunta de agrupación escrita antes de construir la tabla. (2) Filas y Valores configurados con la operación correcta. (3) Resultados ordenados para que el grupo destacado salte a la vista. (4) Captura visible de la tabla dinámica. (5) Conclusión de tres renglones con grupo destacado y decisión posible.
\end{softbox}

\puente{Antes de cerrar, mira la tabla dinámica desde las cinco dimensiones humanas. Agrupar bien es respetar las diferencias. Agrupar mal es disolver las particularidades en un promedio anónimo.}

\milcestacion{milcVino}{Evaluación liberadora · cinco dimensiones}

\begin{softbox}{milcVino}{milcGris}{Sentido de la evaluación}
Evaluar no es solo revisar una nota. Es reconocer qué aprendí, qué cambió en mí, qué emociones manejé, cómo me relaciono con la ciudad y la comunidad, y qué le dejaré a quien viene detrás.
\end{softbox}

\textbf{Mi reflexión liberadora en cinco dimensiones.} Completa cada frase iniciada:

\small
\begin{tabularx}{\textwidth}{>{\bfseries\RaggedRight\arraybackslash}p{3.4cm}Y>{\RaggedRight\arraybackslash}p{4.6cm}}
\rowcolor{milcVino}\color{white}Dimensión & \color{white}\bfseries Mi respuesta (frame starter) & \color{white}\bfseries Algo que descubrí\\
1. Personal & Lo que más cambió en mí fue \linead{6.5cm} & \linead{4.2cm}\\
2. Emocional & La emoción más fuerte fue \linead{2.5cm} y la regulé \linead{3cm} & \linead{4.2cm}\\
3. Ciudadana & En democracia esto importa porque \linead{6cm} & \linead{4.2cm}\\
4. Local (Cartago) & En mi barrio sería útil para \linead{6.5cm} & \linead{4.2cm}\\
5. Intergeneracional & A mi abuela/abuelo le diría \linead{6.5cm} & \linead{4.2cm}\\
\end{tabularx}
\normalsize

\begin{infoband}{milcVino}
\textbf{Para la próxima sustentación.} Vuelve a esta tabla en seis meses. Verás que las respuestas cambian — no porque lo aprendido cambie, sino porque tú cambias. La evaluación liberadora es una conversación contigo a través del tiempo.
\end{infoband}


\puente{Y para terminar, tres voces sobre agrupar. Dussel sobre lo que resiste a toda categoría, Marco Aurelio sobre mirar el conjunto, Floridi sobre los patrones pequeños que solo significan agregados.}

\milcestacion{milcGradeDark}{Triángulo de pensamiento · cierre filosófico}

\begin{softbox}{milcVino}{milcGris}{Dussel · lente del nosotros}
\textit{«El rostro del hombre se revela como otro cuando se recorta en nuestro sistema de instrumentos como exterior, como alguien, como una libertad que interpela, que provoca, que aparece como el que resiste a la totalización instrumental. No es algo; es alguien.»}

{\footnotesize\itshape\color{milcNegro!70}--- Enrique Dussel · Filosofía de la liberación (1977), §2.4.2.2}

Agrupar es meter a alguien en una fila. Cuando agrupas por curso, quien cambió de curso a mitad de año se pierde o se duplica. Cuando agrupas por barrio, quien no tiene dirección fija no entra. Y lo que no encaja en ninguna categoría termina en «otros». Esa fila es la que hay que mirar: ahí está lo que la tabla no supo nombrar.

\textbf{Cuando agrupé mi tabla, ¿quién quedó en «otros» o en «no aplica»? ¿Qué cambiaría si esa fila fuera visible?}
\end{softbox}
\linead{\linewidth}\\[1mm]
\linead{\linewidth}

\begin{softbox}{milcMostaza}{milcGris}{Estoicismo · Marco Aurelio · Meditaciones VII, 47 (c. 175 d.C.) · cuidado interior}
\textit{«Conduce mirar alrededor el curso de los astros, como quien gira con ellos, y contemplar también frecuentemente las mutuas conversiones de los elementos, porque las consideraciones de estas cosas purifican a uno de las manchas de esta vida terrestre.»}

{\footnotesize\itshape\color{milcNegro!70}--- Marco Aurelio · Meditaciones VII, 47 (c. 175 d.C.)}

Subir la mirada al conjunto aclara. La tabla dinámica hace ese movimiento: deja de leer fila por fila y muestra el todo agrupado. Pero conviene bajar también. El promedio del conjunto no dice lo que le pasa a cada grupo, y una cuadrilla entera puede desaparecer dentro de un solo número.

\textbf{¿Qué vi al mirar el conjunto que no se veía fila por fila? ¿Y qué dejé de ver al mirarlo así?}
\end{softbox}
\linead{\linewidth}\\[1mm]
\linead{\linewidth}

\begin{softbox}{milcTurquesa}{milcGris}{Floridi · lente de la infoesfera}
\textit{«Los pequeños patrones solo pueden ser significativos si se agregan correctamente, se comparan y se procesan a tiempo. (trad. propia, abreviada)»}

{\footnotesize\itshape\color{milcNegro!70}--- Luciano Floridi · Big data and their epistemological challenge (2012)}

Una fila sola no dice nada. Una nota, un gasto, una jornada: ruido. Agrupadas por curso, por mes o por tipo, esas mismas filas muestran un patrón. La tabla dinámica es la máquina de agregar correctamente. Pero el «correctamente» es tuyo: la máquina agrupa por donde le digas.

\textbf{¿Mi agrupación deja ver el patrón, o lo fabrica porque agrupé justo por donde quería que saliera?}
\end{softbox}
\linead{\linewidth}\\[1mm]
\linead{\linewidth}

\begin{infoband}{milcNegro}
\textbf{Nota para el docente.} Las tres son citas textuales y llevan su referencia debajo. Puedes remitir a la obra en clase.
\end{infoband}

\milcestacion{milcVino}{Mi autoevaluación · marca una casilla por fila}

{\small
\setlength{\extrarowheight}{2.2pt}
\begin{tabularx}{\textwidth}{>{\RaggedRight\arraybackslash\bfseries}p{3.0cm}YYY}
\rowcolor{milcVino}\color{white}\bfseries Criterio & \color{white}\bfseries Lo logré & \color{white}\bfseries En proceso & \color{white}\bfseries Necesito apoyo\\[.6mm]
Comprensión del tema & \milcopcion{Explico las ideas con mis palabras.} & \milcopcion{Reconozco las ideas, falta precisión.} & \milcopcion{Necesito repasar lo básico.}\\[.8mm]
\rowcolor{milcGris}Pensamiento computacional & \milcopcion{Usé los pilares al armar mi producto.} & \milcopcion{Usé algunos pilares.} & \milcopcion{Necesito releer la estación 2.}\\[.8mm]
Aplicación y producto & \milcopcion{Mi producto cumple los criterios.} & \milcopcion{Está armado, con ajustes pendientes.} & \milcopcion{Aún está en construcción.}\\[.8mm]
\rowcolor{milcGris}Reflexión 5 dimensiones & \milcopcion{Respondí cada una con honestidad.} & \milcopcion{Respondí algunas con detalle.} & \milcopcion{Mis respuestas son muy generales.}\\[.8mm]
Triángulo de pensamiento & \milcopcion{Conecté las 3 voces con mi trabajo.} & \milcopcion{Conecté al menos una.} & \milcopcion{No las relaciono con mi proceso.}\\[.8mm]
\end{tabularx}}

\vspace{3mm}
\textbf{Hoy entendí que:}\\[1mm]
\linead{\linewidth}\\[1mm]
\linead{\linewidth}

\textbf{Mi compromiso para la próxima sesión es:}\\[1mm]
\linead{\linewidth}\\[1mm]
\linead{\linewidth}

\begin{softbox}{milcMostaza}{milcGris}{Cita de despedida · Marco Aurelio}
\textit{``El obstáculo en el camino se vuelve el camino. Lo que estorba al camino se vuelve camino.''}\\[1mm]
Cada error de hoy, bien observado, se convierte en pieza del camino que pisarás mañana.
\end{softbox}

\small
\begin{tabularx}{\textwidth}{>{\bfseries}p{3.6cm}Y}
\rowcolor{milcGradeDark!12}Nombre completo: & \linead{10cm}\\
Fecha: & \linead{4cm} \quad Curso/grupo: \linead{4cm}\\
Mi firma: & \linead{9cm}\\
\end{tabularx}
\normalsize

\begin{infoband}{milcGradeDark}
\textbf{Para la próxima semana.} Aplica una tabla dinámica real a una tabla cotidiana fuera del aula: gastos por categoría, calificaciones por área, horas de sueño por día de la semana. Documenta el grupo destacado, el grupo en riesgo y una decisión real que tomes a partir de esa comparación. La phronesis no se queda en la hoja: se entrena en lo que decides después.
\end{infoband}

\end{document}
